% =============================================================================
% paper46_galois_theory_en.tex
% Paper 46: Galois Theory — From Field Extensions to the Insolvability of the Quintic
% Specialist Mathematics Project — Build 135
% Author: Michael Fuhrmann
% Date: 2026-03-12
% =============================================================================

\documentclass[12pt,a4paper]{amsart}
\usepackage{amsmath,amssymb,amsthm}
\usepackage{mathtools}
\usepackage[utf8]{inputenc}
\usepackage[T1]{fontenc}
\usepackage{hyperref}
\usepackage{enumitem,booktabs,array}
\usepackage{geometry}
\geometry{margin=2.5cm}

% ---------------------------------------------------------------------------
% Theorem environments
% ---------------------------------------------------------------------------
\newtheorem{theorem}{Theorem}[section]
\newtheorem{lemma}[theorem]{Lemma}
\newtheorem{corollary}[theorem]{Corollary}
\newtheorem{proposition}[theorem]{Proposition}
\newtheorem{conjecture}[theorem]{Conjecture}
\theoremstyle{definition}
\newtheorem{definition}[theorem]{Definition}
\newtheorem{example}[theorem]{Example}
\newtheorem{remark}[theorem]{Remark}

% ---------------------------------------------------------------------------
% Metadata
% ---------------------------------------------------------------------------
\title{Galois Theory:\\
       From Field Extensions to the Insolvability of the Quintic}

\author{Michael Fuhrmann}
\address{Specialist Mathematics Project, Build 135}
\email{specialist-maths@localhost}
\date{\today}

\keywords{Galois theory, field extensions, splitting fields, Galois group,
          fundamental theorem, solvability by radicals, Abel-Ruffini theorem,
          Kronecker-Weber theorem, constructibility}

\subjclass[2020]{12F10, 12F20, 11R32, 12E05}

% ---------------------------------------------------------------------------
% Macros
% ---------------------------------------------------------------------------
\newcommand{\Q}{\mathbb{Q}}
\newcommand{\R}{\mathbb{R}}
\newcommand{\C}{\mathbb{C}}
\newcommand{\Z}{\mathbb{Z}}
\newcommand{\F}{\mathbb{F}}
\newcommand{\Gal}{\operatorname{Gal}}
\newcommand{\Aut}{\operatorname{Aut}}
\newcommand{\ord}{\operatorname{ord}}
\newcommand{\lcm}{\operatorname{lcm}}
\newcommand{\Fix}{\operatorname{Fix}}
\newcommand{\chara}{\operatorname{char}}

% =============================================================================
\begin{document}
% =============================================================================

\begin{abstract}
Galois theory is one of the most profound achievements in the history of
mathematics. Developed by \'{E}variste Galois in 1832 and completed
posthumously, it establishes a deep correspondence between the symmetry groups
of polynomial roots and the arithmetic of field extensions. This paper presents
a self-contained introduction to Galois theory, starting from the foundations
of field extensions and algebraic closures, proceeding through splitting fields,
separability, normality, and the Galois group, to the celebrated Fundamental
Theorem of Galois Theory. We then apply this machinery to settle the classical
problems of solvability by radicals—proving the Abel-Ruffini Theorem—and
constructibility by compass and straightedge. The Kronecker-Weber Theorem on
abelian extensions of $\Q$ is presented as a further landmark application.
Throughout, concrete examples illustrate the abstract theory.
\end{abstract}

\maketitle
\tableofcontents

% =============================================================================
\section{Introduction}
% =============================================================================

The problem of solving polynomial equations by radicals—\emph{closed-form
formulas} involving the coefficients and nested radicals—occupied mathematicians
for centuries. The Babylonians knew the quadratic formula around 2000 BCE;
Cardano published the cubic formula in \emph{Ars Magna} (1545), and Ferrari
derived a formula for degree four in the same decade. But degree five resisted
all attacks until 1799, when Paolo Ruffini announced (with an incomplete proof)
that no such formula could exist; Abel gave the first rigorous proof in 1824.

The definitive explanation came from \'{E}variste Galois (1811--1832), who in
a memoir written the night before his fatal duel created what we now call
\emph{Galois theory}. Galois associated to every polynomial $f \in K[x]$ a
group $\Gal(f/K)$ of symmetries of its roots, and proved:
\[
  f \text{ is solvable by radicals} \iff \Gal(f/K) \text{ is a solvable group.}
\]
Since the generic degree-five polynomial has Galois group $S_5$, and $A_5
\trianglelefteq S_5$ is simple and non-abelian, $S_5$ is not solvable, whence
no radical formula exists for degree five.

The theory has grown far beyond its origin. Today Galois theory is central to
algebraic number theory, the Langlands programme, and modern cryptography.

\subsection*{Historical Timeline}

\begin{center}
\begin{tabular}{ll}
\toprule
Year & Development \\
\midrule
$\sim$2000 BCE & Babylonian quadratic formula \\
1545 & Cardano/Ferrari: cubic and quartic formulas (\emph{Ars Magna}) \\
1799 & Ruffini: first (incomplete) proof that degree 5 is unsolvable \\
1824 & Abel: rigorous proof of Abel-Ruffini \\
1832 & Galois: group-theoretic criterion; memoir published posthumously \\
1870 & Jordan: first book on group theory including Galois theory \\
1801/1853 & Gauss/Kronecker-Weber: cyclotomic fields and abelian extensions \\
\bottomrule
\end{tabular}
\end{center}

% =============================================================================
\section{Field Extensions}
\label{sec:fieldext}
% =============================================================================

\subsection{Basic Definitions}

\begin{definition}
A \emph{field extension} $L/K$ is an inclusion of fields $K \subseteq L$.
The \emph{degree} $[L:K] = \dim_K L$ is the dimension of $L$ as a
$K$-vector space. We say $L/K$ is \emph{finite} if $[L:K] < \infty$.
\end{definition}

\begin{theorem}[Tower Law]
If $K \subseteq F \subseteq L$ are fields, then
\[
  [L:K] = [L:F] \cdot [F:K].
\]
\end{theorem}

\begin{proof}
Let $\{e_i\}$ be a basis of $F$ over $K$ and $\{f_j\}$ a basis of $L$ over
$F$. One verifies that $\{e_i f_j\}$ is a $K$-basis of $L$, giving
$[L:K] = |I| \cdot |J| = [F:K] \cdot [L:F]$.
\end{proof}

\subsection{Algebraic and Transcendental Elements}

\begin{definition}
An element $\alpha \in L$ is \emph{algebraic over $K$} if there exists a
nonzero polynomial $f \in K[x]$ with $f(\alpha) = 0$. Otherwise $\alpha$ is
\emph{transcendental over $K$}. The extension $L/K$ is \emph{algebraic} if
every element of $L$ is algebraic over $K$.
\end{definition}

\begin{definition}
The \emph{minimal polynomial} $\mathrm{min}_{K}(\alpha)$ of an algebraic
element $\alpha$ over $K$ is the unique monic irreducible polynomial in $K[x]$
that vanishes at $\alpha$.
\end{definition}

\begin{proposition}
If $\alpha$ is algebraic over $K$ with minimal polynomial $p$ of degree $n$,
then $K(\alpha) \cong K[x]/(p)$ and $[K(\alpha):K] = n$.
\end{proposition}

\begin{example}
The minimal polynomial of $\sqrt{2}$ over $\Q$ is $x^2 - 2$, so
$[\Q(\sqrt{2}):\Q] = 2$ and $\Q(\sqrt{2}) = \{a + b\sqrt{2} : a,b \in \Q\}$.
Similarly, $\mathrm{min}_\Q(\sqrt[3]{2}) = x^3 - 2$, giving
$[\Q(\sqrt[3]{2}):\Q] = 3$.
\end{example}

\begin{example}
The algebraic closure $\overline{\Q}$ of $\Q$ is algebraic over $\Q$ but of
infinite degree. The real numbers $\pi$ and $e$ are transcendental over $\Q$
(Lindemann--Weierstrass).
\end{example}

\subsection{Simple and Compositum Extensions}

\begin{definition}
The \emph{compositum} $L_1 L_2$ of two extensions $L_1, L_2$ of $K$ inside
a common overfield is the smallest subfield containing both $L_1$ and $L_2$.
\end{definition}

% =============================================================================
\section{Splitting Fields and Algebraic Closure}
% =============================================================================

\subsection{Splitting Fields}

\begin{definition}
A \emph{splitting field} of $f \in K[x]$ is a field extension $L/K$ such that
$f$ splits into linear factors over $L$, i.e.\ $f(x) = a \prod_{i=1}^n (x -
\alpha_i)$ with $\alpha_i \in L$, and $L = K(\alpha_1,\dots,\alpha_n)$.
\end{definition}

\begin{theorem}[Existence and Uniqueness of Splitting Fields]
\label{thm:splitfield}
For every $f \in K[x]$ there exists a splitting field $L/K$. It is unique up
to $K$-isomorphism, and $[L:K] \leq (\deg f)!$
\end{theorem}

\begin{proof}[Proof sketch]
\emph{Existence}: By induction on $\deg f$. If $f$ has no irreducible factor of
degree $> 1$ in $K[x]$, then $K$ itself is the splitting field. Otherwise pick
an irreducible factor $p \mid f$, set $K_1 = K[x]/(p)$, adjoin a root
$\alpha_1$ of $p$, and continue with $f/(x-\alpha_1)$ over $K_1$.

\emph{Uniqueness}: Any two splitting fields $L, L'$ are connected by a
$K$-isomorphism, constructed by extending isomorphisms step by step along the
roots.
\end{proof}

\begin{example}
The splitting field of $x^3 - 2$ over $\Q$ is
$\Q(\sqrt[3]{2}, \omega)$ where $\omega = e^{2\pi i/3}$ is a primitive cube
root of unity. We have $[\Q(\sqrt[3]{2},\omega):\Q] = 6$.
\end{example}

\subsection{Algebraic Closure}

\begin{definition}
A field $\Omega$ is \emph{algebraically closed} if every non-constant polynomial
in $\Omega[x]$ has a root in $\Omega$. The \emph{algebraic closure}
$\overline{K}$ of $K$ is the smallest algebraically closed field containing $K$.
\end{definition}

\begin{theorem}[Existence of Algebraic Closure (Steinitz, 1910)]
Every field $K$ has an algebraic closure $\overline{K}$, unique up to
$K$-isomorphism.
\end{theorem}

% =============================================================================
\section{Separability and Normality}
% =============================================================================

\subsection{Separable Extensions}

\begin{definition}
A polynomial $f \in K[x]$ is \emph{separable} if it has no repeated roots in
$\overline{K}$, equivalently if $\gcd(f, f') = 1$ in $K[x]$. An algebraic
element $\alpha$ over $K$ is separable if $\mathrm{min}_K(\alpha)$ is
separable. An algebraic extension $L/K$ is \emph{separable} if every element
of $L$ is separable over $K$.
\end{definition}

\begin{remark}
Every algebraic extension of a field of characteristic $0$ (e.g.\ $\Q$) or of
a finite field is separable. Non-separable extensions arise only in
characteristic $p > 0$.
\end{remark}

\begin{theorem}[Primitive Element Theorem]
If $L/K$ is a finite separable extension, then $L = K(\theta)$ for some
$\theta \in L$.
\end{theorem}

\subsection{Normal Extensions}

\begin{definition}
A finite extension $L/K$ is \emph{normal} if it is the splitting field of some
polynomial $f \in K[x]$. Equivalently, for every irreducible $p \in K[x]$:
if $p$ has one root in $L$, then $p$ splits completely in $L$.
\end{definition}

\begin{example}
$\Q(\sqrt[3]{2})/\Q$ is \emph{not} normal: $x^3 - 2$ has one root
$\sqrt[3]{2} \in \Q(\sqrt[3]{2})$ but does not split completely there.
Its splitting field $\Q(\sqrt[3]{2},\omega)$ is normal over $\Q$.
\end{example}

% =============================================================================
\section{The Galois Group}
% =============================================================================

\begin{definition}
Let $L/K$ be a field extension. The \emph{Galois group} is
\[
  \Gal(L/K) = \Aut_K(L) = \{\sigma: L \xrightarrow{\sim} L \mid \sigma|_K = \mathrm{id}_K\}.
\]
This is a group under composition.
\end{definition}

\begin{definition}
A finite extension $L/K$ is a \emph{Galois extension} if it is both normal and
separable. In this case $|\Gal(L/K)| = [L:K]$.
\end{definition}

\begin{example}[Cyclotomic Extension]
Let $\zeta_n = e^{2\pi i/n}$. The extension $\Q(\zeta_n)/\Q$ is Galois with
\[
  \Gal(\Q(\zeta_n)/\Q) \cong (\Z/n\Z)^\times.
\]
Each automorphism $\sigma_k$ (with $\gcd(k,n)=1$) sends $\zeta_n \mapsto \zeta_n^k$.
\end{example}

\begin{example}[Galois Group of $x^3 - 2$]
The splitting field $L = \Q(\sqrt[3]{2}, \omega)$ over $\Q$ has $[L:\Q] = 6$.
The Galois group $\Gal(L/\Q) \cong S_3$, the symmetric group on three elements,
acting by permuting the three roots $\{\sqrt[3]{2}, \omega\sqrt[3]{2},
\omega^2\sqrt[3]{2}\}$.
\end{example}

\begin{example}[Galois Group of a Quadratic]
For $f(x) = x^2 - d$ with $d \in \Q$ not a perfect square,
$\Gal(\Q(\sqrt{d})/\Q) \cong \Z/2\Z$,
generated by $\sqrt{d} \mapsto -\sqrt{d}$.
\end{example}

\subsection{Discriminant and Galois Group}

\begin{theorem}
Let $f \in K[x]$ be separable of degree $n$, with splitting field $L/K$.
The Galois group $G = \Gal(L/K)$ embeds in $S_n$ via its action on the roots.
If $\Delta = \prod_{i < j}(\alpha_i - \alpha_j)^2$ is the discriminant:
\begin{enumerate}[label=(\roman*)]
  \item $G \subseteq A_n$ if and only if $\sqrt{\Delta} \in K$.
  \item For $n = 3$: $G \cong A_3 \cong \Z/3\Z$ iff $\sqrt{\Delta} \in K$;
        otherwise $G \cong S_3$.
  \item For $n = 2$: $G \cong \Z/2\Z$ iff $\Delta \notin K^{\times 2}$.
\end{enumerate}
\end{theorem}

\subsection{Frobenius Endomorphism}

\begin{definition}
For a finite field $\F_{p^n}$, the \emph{Frobenius automorphism}
is $\varphi: x \mapsto x^p$. Then
\[
  \Gal(\F_{p^n}/\F_p) = \langle \varphi \rangle \cong \Z/n\Z.
\]
\end{definition}

% =============================================================================
\section{The Fundamental Theorem of Galois Theory}
% =============================================================================

\begin{theorem}[Fundamental Theorem of Galois Theory]
\label{thm:ftgt}
Let $L/K$ be a finite Galois extension with Galois group $G = \Gal(L/K)$.
There is an order-reversing bijection
\[
  \left\{\text{subgroups } H \leq G\right\}
  \longleftrightarrow
  \left\{\text{intermediate fields } K \subseteq F \subseteq L\right\}
\]
given by $H \mapsto L^H = \Fix(H) = \{x \in L : \sigma(x) = x\ \forall \sigma \in H\}$
and $F \mapsto \Gal(L/F)$. Moreover:
\begin{enumerate}[label=(\roman*)]
  \item $[L:L^H] = |H|$ and $[L^H:K] = [G:H]$.
  \item $H$ is normal in $G$ if and only if $L^H/K$ is a Galois extension, in
        which case $\Gal(L^H/K) \cong G/H$.
  \item The trivial subgroup $\{e\}$ corresponds to $L$, and $G$ itself
        corresponds to $K$.
\end{enumerate}
\end{theorem}

\begin{proof}[Proof sketch]
The key tools are:
\begin{itemize}
  \item \emph{Artin's theorem}: If $G$ is a finite group of automorphisms of $L$
        with fixed field $K = L^G$, then $[L:K] = |G|$ and $L/K$ is Galois with
        group $G$.
  \item \emph{Injectivity}: The map $H \mapsto L^H$ is injective because any
        $\sigma \in H \setminus \{e\}$ moves some element of $L$ not fixed by all
        of $H$.
  \item \emph{Surjectivity}: For any intermediate field $F$, we have
        $F = L^{\Gal(L/F)}$ by Artin applied to $\Gal(L/F)$ acting on $L$ with
        fixed field containing $F$.
  \item The order-reversal $[L:L^H] = |H|$ follows directly from Artin's theorem.
\end{itemize}
\end{proof}

\begin{example}[Galois Correspondence for $\Q(\sqrt{2},\sqrt{3})$]
Let $L = \Q(\sqrt{2},\sqrt{3})$. Then $[L:\Q] = 4$ and
$G = \Gal(L/\Q) \cong \Z/2\Z \times \Z/2\Z = \{e, \sigma, \tau, \sigma\tau\}$
where $\sigma(\sqrt{2}) = -\sqrt{2}$, $\sigma(\sqrt{3}) = \sqrt{3}$,
$\tau(\sqrt{2}) = \sqrt{2}$, $\tau(\sqrt{3}) = -\sqrt{3}$.

The subgroup lattice and corresponding field lattice:
\begin{center}
\begin{tabular}{ll}
\toprule
Subgroup $H$ & Fixed field $L^H$ \\
\midrule
$\{e\}$ & $L = \Q(\sqrt{2},\sqrt{3})$ \\
$\{e, \sigma\}$ & $\Q(\sqrt{3})$ \\
$\{e, \tau\}$ & $\Q(\sqrt{2})$ \\
$\{e, \sigma\tau\}$ & $\Q(\sqrt{6})$ \\
$G$ & $\Q$ \\
\bottomrule
\end{tabular}
\end{center}
All subgroups are normal (since $G$ is abelian), corresponding to the fact that
all intermediate extensions are Galois over $\Q$.
\end{example}

% =============================================================================
\section{Solvability by Radicals}
% =============================================================================

\subsection{Radical Extensions}

\begin{definition}
An extension $L/K$ is a \emph{radical extension} if there exists a tower
\[
  K = K_0 \subset K_1 \subset \cdots \subset K_r = L
\]
with $K_{i+1} = K_i(\alpha_i)$ and $\alpha_i^{n_i} \in K_i$ for some
$n_i \in \Z_{>0}$.
A polynomial $f \in K[x]$ is \emph{solvable by radicals} if its splitting
field is contained in a radical extension of $K$.
\end{definition}

\begin{definition}
A group $G$ is \emph{solvable} if it has a subnormal series
\[
  \{e\} = G_r \trianglelefteq G_{r-1} \trianglelefteq \cdots \trianglelefteq G_0 = G
\]
with each factor $G_{i-1}/G_i$ abelian.
\end{definition}

\begin{example}[Solvable Groups of Small Order]
$S_1, S_2, S_3, S_4$ are solvable:
\begin{align*}
  S_3 &\supset A_3 \cong \Z/3\Z \supset \{e\},
  \quad \text{factors: } \Z/2\Z,\ \Z/3\Z \\
  S_4 &\supset A_4 \supset V_4 \cong (\Z/2\Z)^2 \supset \Z/2\Z \supset \{e\},
  \quad \text{all factors abelian.}
\end{align*}
\end{example}

\subsection{The Main Theorem on Radical Solvability}

\begin{theorem}[Galois, 1832]
\label{thm:solvable}
Let $K$ be a field of characteristic $0$ and $f \in K[x]$. Then $f$ is
solvable by radicals over $K$ if and only if $\Gal(f/K)$ is a solvable group.
\end{theorem}

\begin{proof}[Proof sketch]
($\Rightarrow$) If $L$ lies in a radical tower $K = K_0 \subset \cdots \subset
K_r$, one can replace each step by a cyclic extension (after adjoining roots of
unity), yielding a chain of Galois groups with cyclic---hence abelian---quotients,
so $\Gal(L/K)$ is solvable.

($\Leftarrow$) If $G = \Gal(L/K)$ is solvable with series
$\{e\} = G_r \trianglelefteq \cdots \trianglelefteq G_0 = G$, the fixed fields
$K_i = L^{G_i}$ form intermediate extensions with
$\Gal(K_{i+1}/K_i) \cong G_i/G_{i+1}$ cyclic. By Kummer theory (after
adjoining appropriate roots of unity), each cyclic extension of degree $n$ is
of the form $K_i(\alpha_i)$ with $\alpha_i^n \in K_i$, thus radical.
\end{proof}

\subsection{Explicit Formulas for Degrees 2, 3, and 4}

For completeness we record the classical radical formulas.

\begin{example}[Degree 2]
$x^2 + px + q = 0$ has solutions
$x = \dfrac{-p \pm \sqrt{p^2 - 4q}}{2}$.
The Galois group is trivial if $p^2 - 4q$ is a square in $K$, and $\Z/2\Z$
otherwise; both are solvable.
\end{example}

\begin{example}[Degree 3: Cardano's Formula]
For $x^3 + px + q = 0$ (depressed form), set $\Delta = -4p^3 - 27q^2$.
The roots are
\[
  x_k = \omega^k \sqrt[3]{-\frac{q}{2} + \sqrt{-\frac{\Delta}{108}}}
       + \omega^{-k} \sqrt[3]{-\frac{q}{2} - \sqrt{-\frac{\Delta}{108}}},
  \quad k = 0,1,2,
\]
where $\omega = e^{2\pi i/3}$. The radical tower is
$\Q \subset \Q(\omega) \subset \Q(\omega, u)$ where $u$ is a cube root.
\end{example}

\begin{example}[Degree 4: Ferrari's Method]
Ferrari (1545) reduced the quartic to a cubic (\emph{resolvent}) via an
auxiliary substitution. If $f(x) = x^4 + bx^2 + cx + d$, the resolvent cubic
is $y^3 - by^2 - 4dy + (4bd - c^2)$. Solving the cubic by Cardano, one
reduces to two quadratics, giving a radical tower of depth 4.
\end{example}

% =============================================================================
\section{The Insolvability of the Quintic}
% =============================================================================

\subsection{Simplicity of $A_5$}

\begin{theorem}[$A_5$ is Simple]
\label{thm:a5simple}
The alternating group $A_5$ of even permutations of $5$ elements is simple,
i.e.\ it has no proper non-trivial normal subgroup.
\end{theorem}

\begin{proof}
$|A_5| = 60$. The conjugacy classes in $A_5$ have sizes
$1, 12, 12, 15, 20$. A normal subgroup $N \trianglelefteq A_5$ is a union of
conjugacy classes containing the identity; so $|N|$ divides $60$ and must be
a sum from $\{1, 12, 12, 15, 20\}$ including $1$. The only possible sums
dividing $60$ are $1$ and $60$, so $N = \{e\}$ or $N = A_5$.
\end{proof}

\subsection{Insolvability of $S_5$}

\begin{theorem}[$S_5$ is not solvable]
\label{thm:s5notsolv}
The symmetric group $S_5$ is not solvable.
\end{theorem}

\begin{proof}
The only normal subgroups of $S_5$ are $\{e\}$, $A_5$, and $S_5$. The series
$S_5 \supset A_5 \supset \{e\}$ has factor $A_5/\{e\} \cong A_5$, which is not
abelian ($|A_5| = 60$, and simple non-trivial groups of order $> 1$ are not
abelian unless cyclic). Hence $S_5$ is not solvable.
\end{proof}

\subsection{The Abel-Ruffini-Galois Theorem}

\begin{theorem}[Abel-Ruffini, 1824; Galois, 1832]
\label{thm:abelruffini}
There is no general formula by radicals for the roots of a polynomial equation
of degree $n \geq 5$ over $\Q$.
\end{theorem}

\begin{proof}
It suffices to exhibit one polynomial of degree $5$ whose Galois group is $S_5$.
Consider $f(x) = x^5 - x - 1 \in \Q[x]$. One can verify (using the
discriminant and Dedekind's theorem on reduction modulo primes) that
$\Gal(f/\Q) \cong S_5$. By Theorem~\ref{thm:solvable}, $f$ is solvable by
radicals if and only if $S_5$ is solvable. But $S_5$ is not solvable
(Theorem~\ref{thm:s5notsolv}). Hence $f$ has no expression of its roots in
terms of radicals and the coefficients.

Since the \emph{generic} degree-$n$ polynomial over $\Q(t_1,\dots,t_n)$ has
Galois group $S_n$, and $S_n$ is not solvable for $n \geq 5$, no general
radical formula exists for any degree $n \geq 5$.
\end{proof}

\begin{remark}
The theorem does \emph{not} assert that no degree-5 polynomial has roots
expressible by radicals. Many do: for instance $x^5 - 2$ is solvable
(its Galois group is the Frobenius group $\Z/5\Z \rtimes \Z/4\Z$ of order 20,
which is solvable). The assertion is that no \emph{universal formula} in the
coefficients works.
\end{remark}

\begin{remark}
Composition series by degree:
\begin{center}
\begin{tabular}{lccc}
\toprule
Degree & Galois group & Solvable? & Formula \\
\midrule
1 & $\{e\}$ & Yes & $x = -b/a$ \\
2 & $\Z/2\Z$ & Yes & Quadratic formula \\
3 & $\subseteq S_3$ & Yes & Cardano (1545) \\
4 & $\subseteq S_4$ & Yes & Ferrari (1545) \\
$\geq 5$ & can be $S_n$ & No & None in general \\
\bottomrule
\end{tabular}
\end{center}
\end{remark}

% =============================================================================
\section{The Kronecker-Weber Theorem}
% =============================================================================

\subsection{Cyclotomic Fields}

\begin{definition}
The $n$-th \emph{cyclotomic polynomial} is
\[
  \Phi_n(x) = \prod_{\substack{1 \leq k \leq n \\ \gcd(k,n) = 1}} (x - e^{2\pi i k/n})
  \in \Z[x],
\]
of degree $\varphi(n) = |(\Z/n\Z)^\times|$. The $n$-th cyclotomic field is
$\Q(\zeta_n)$ where $\zeta_n = e^{2\pi i/n}$.
\end{definition}

\begin{theorem}[Cyclotomic Galois Group]
$\Q(\zeta_n)/\Q$ is a Galois extension of degree $\varphi(n)$, with
$\Gal(\Q(\zeta_n)/\Q) \cong (\Z/n\Z)^\times$
(abelian group).
\end{theorem}

\subsection{The Theorem}

\begin{theorem}[Kronecker-Weber, 1853/1886]
\label{thm:kronweber}
Every finite abelian extension of $\Q$ is contained in a cyclotomic field.
More precisely, if $L/\Q$ is a finite Galois extension with abelian Galois
group, then $L \subseteq \Q(\zeta_n)$ for some $n \in \Z_{>0}$.
\end{theorem}

\begin{remark}
Kronecker stated the theorem in 1853 but without complete proof. Weber gave
an (almost complete) proof in 1886. The first complete proof is due to
Hilbert (1896).
\end{remark}

\begin{proof}[Proof sketch via Local-Global Arguments]
The modern proof proceeds in two stages:
\begin{enumerate}
  \item \emph{Local Kronecker-Weber}: Every finite abelian extension of $\Q_p$
        (the $p$-adic rationals) is contained in a cyclotomic field over $\Q_p$
        (proved using the structure of $\Q_p^\times$ via local class field theory
        or direct arguments).
  \item \emph{Global assembly}: A finite abelian extension $L/\Q$ is determined
        by its completions at all primes. Using the local result at each prime
        $p$ dividing $[L:\Q]$, one constructs an $n$ such that $L \subseteq
        \Q(\zeta_n)$.
\end{enumerate}
\end{proof}

\begin{example}
The extension $\Q(\sqrt{d})/\Q$ for $d$ squarefree has Galois group $\Z/2\Z$
(abelian). By Kronecker-Weber, $\Q(\sqrt{d}) \subseteq \Q(\zeta_n)$ for some
$n$. Indeed: $\Q(\sqrt{d}) \subseteq \Q(\zeta_{|d|})$ (or $\Q(\zeta_{4|d|}$
for $d \equiv 2,3 \pmod{4}$).
\end{example}

\begin{example}
$\Q(i) = \Q(\sqrt{-1}) \subseteq \Q(\zeta_4)$ since $\zeta_4 = i$.
\end{example}

\begin{remark}
Kronecker-Weber is the first major theorem of class field theory. Its
generalization—the Artin reciprocity law—describes abelian extensions of any
number field $K$ in terms of the idèle class group of $K$.
\end{remark}

% =============================================================================
\section{Constructibility by Compass and Straightedge}
% =============================================================================

\subsection{Algebraic Criterion}

\begin{definition}
A real number $\alpha$ is \emph{constructible} if it can be obtained from $0$
and $1$ by a finite sequence of field operations and square roots. Equivalently,
$\alpha$ is constructible if and only if there exists a tower
\[
  \Q = F_0 \subset F_1 \subset \cdots \subset F_m
\]
with $[F_{i+1}:F_i] = 2$ for each $i$ and $\alpha \in F_m$.
\end{definition}

\begin{theorem}[Algebraic Criterion for Constructibility]
\label{thm:constructible}
A real number $\alpha$ is constructible if and only if $[\Q(\alpha):\Q]$ is a
power of $2$.
\end{theorem}

\subsection{Classical Impossibility Results}

\begin{corollary}[Doubling the Cube — Impossible]
It is impossible to construct $\sqrt[3]{2}$ by compass and straightedge.
\end{corollary}

\begin{proof}
$[\Q(\sqrt[3]{2}):\Q] = 3$, which is not a power of $2$.
\end{proof}

\begin{corollary}[Trisecting an Angle — Impossible in General]
A general angle of $60°$ cannot be trisected by compass and straightedge.
\end{corollary}

\begin{proof}
Trisecting $60°$ requires constructing $\cos 20°$, which satisfies
$8x^3 - 6x - 1 = 0$ (irreducible over $\Q$, degree 3). Hence
$[\Q(\cos 20°):\Q] = 3$, not a power of $2$.
\end{proof}

\begin{corollary}[Squaring the Circle — Impossible]
The real number $\sqrt{\pi}$ is transcendental (Lindemann, 1882), hence not
algebraic, and in particular not constructible.
\end{corollary}

\subsection{Gauss-Wantzel Theorem: Regular Polygons}

\begin{definition}
The \emph{Fermat primes} are primes of the form $F_k = 2^{2^k} + 1$.
The known Fermat primes are $3, 5, 17, 257, 65537$ (for $k = 0,1,2,3,4$).
\end{definition}

\begin{theorem}[Gauss-Wantzel, 1801/1837]
\label{thm:gausswantzel}
A regular $n$-gon is constructible by compass and straightedge if and only if
\[
  n = 2^a \cdot p_1 \cdot p_2 \cdots p_k,
\]
where $a \geq 0$ and $p_1 < p_2 < \cdots < p_k$ are distinct Fermat primes.
\end{theorem}

\begin{proof}[Proof sketch]
Constructing a regular $n$-gon is equivalent to constructing $\zeta_n = e^{2\pi i/n}$,
which lies in $\Q(\zeta_n)$ with $[\Q(\zeta_n):\Q] = \varphi(n)$.
By Theorem~\ref{thm:constructible}, this is possible iff $\varphi(n)$ is a
power of $2$. Elementary number theory shows $\varphi(n) = 2^a$ iff $n$ has
the stated form.
\end{proof}

\begin{example}[Heptadecagon]
Gauss (age 19, 1796) discovered that the regular 17-gon ($n = 17 = F_2$)
is constructible, since $\varphi(17) = 16 = 2^4$. The explicit construction
involves $\cos(2\pi/17)$, which can be expressed as
\[
  \cos\!\left(\frac{2\pi}{17}\right)
  = \frac{-1 + \sqrt{17} + \sqrt{34 - 2\sqrt{17}}
    + 2\sqrt{17 + 3\sqrt{17} - \sqrt{34-2\sqrt{17}} - 2\sqrt{34+2\sqrt{17}}}}{16}.
\]
This was so surprising to Gauss that it reportedly inspired him to pursue
mathematics rather than philology.
\end{example}

\begin{example}
The regular heptagon ($n = 7$) is \emph{not} constructible: $\varphi(7) = 6$
is not a power of $2$.
\end{example}

% =============================================================================
\section{Summary and Connections}
% =============================================================================

Galois theory ties together algebra, geometry, and arithmetic through a single
unifying principle: \emph{symmetry governs solvability}. The main results
established here form a coherent arc:

\begin{enumerate}
  \item Field extensions of degree $n$ correspond to polynomials of degree $n$
        (§\ref{sec:fieldext} -- §\ref{sec:splitting}).
  \item Galois extensions have a symmetry group (§\ref{sec:galoisgroup}).
  \item The Fundamental Theorem creates a perfect duality between subgroups
        and intermediate fields (Theorem~\ref{thm:ftgt}).
  \item A polynomial is solvable by radicals iff its symmetry group is solvable
        (Theorem~\ref{thm:solvable}).
  \item $S_5$ is not solvable, proving the Abel-Ruffini Theorem
        (Theorem~\ref{thm:abelruffini}).
  \item Abelian extensions of $\Q$ are cyclotomic (Theorem~\ref{thm:kronweber}).
  \item Constructibility reduces to degree being a power of 2
        (Theorem~\ref{thm:constructible}, \ref{thm:gausswantzel}).
\end{enumerate}

\noindent\textbf{Further Directions.}
Modern extensions of Galois theory include:
\begin{itemize}
  \item \emph{Inverse Galois Problem}: Is every finite group realizable as a
        Galois group over $\Q$? (Open in general; solved for solvable groups
        by Shafarevich.)
  \item \emph{\'Etale cohomology and Grothendieck}: Generalisation to algebraic
        geometry via the fundamental group $\pi_1^{\text{\'et}}(\mathrm{Spec}\,K)
        = \Gal(\overline{K}/K)$.
  \item \emph{Langlands Programme}: Deep conjectures relating Galois
        representations $\rho: \Gal(\overline{\Q}/\Q) \to GL_n(\C)$ to
        automorphic forms.
  \item \emph{$p$-adic and Iwasawa Theory}: Galois groups of infinite towers of
        number fields.
\end{itemize}

\addtocounter{section}{-1}
\refstepcounter{section}
\renewcommand{\thesection}{\Alph{section}}

% =============================================================================
\begin{thebibliography}{99}
% =============================================================================

\bibitem{galois1832}
\'{E}.~Galois,
\textit{M\'emoire sur les conditions de r\'esolubilit\'e des \'equations par
radicaux},
Manuscript (1832); first published in \textit{Journal de Math\'ematiques
Pures et Appliqu\'ees} \textbf{11} (1846), 381--444.

\bibitem{abel1824}
N.H.~Abel,
\textit{M\'emoire sur les \'equations alg\'ebriques o\`{u} on d\'emontre
l'impossibilit\'e de la r\'esolution de l'\'equation g\'en\'erale du cinqui\`eme
degr\'e},
Christiania, 1824.

\bibitem{ruffini1799}
P.~Ruffini,
\textit{Teoria generale delle equazioni},
Bologna, 1799.

\bibitem{lang2002}
S.~Lang,
\textit{Algebra}, 3rd ed.,
Graduate Texts in Mathematics \textbf{211},
Springer, New York, 2002.

\bibitem{neukirch1999}
J.~Neukirch,
\textit{Algebraic Number Theory},
Grundlehren der mathematischen Wissenschaften \textbf{322},
Springer, Berlin, 1999.

\bibitem{stewart2015}
I.~Stewart,
\textit{Galois Theory}, 4th ed.,
CRC Press, Boca Raton, 2015.

\bibitem{dummitfoote2004}
D.S.~Dummit and R.M.~Foote,
\textit{Abstract Algebra}, 3rd ed.,
Wiley, Hoboken, 2004.

\bibitem{cox2004}
D.A.~Cox,
\textit{Galois Theory},
Wiley-Interscience, Hoboken, 2004.

\bibitem{artin1944}
E.~Artin,
\textit{Galois Theory},
Notre Dame Mathematical Lectures \textbf{2},
University of Notre Dame Press, 1944.

\bibitem{serre1979}
J.-P.~Serre,
\textit{Local Fields},
Graduate Texts in Mathematics \textbf{67},
Springer, New York, 1979.

\end{thebibliography}

\end{document}
